\begin{textsample}{POS Dim 3 – am – Score 49.00 – am/am\_inf\_21.txt}  \label{ex:f3_pos_031}
8.
Os direitos de aprendizagem e desenvolvimento Ao pensarmos o currículo da Educação \textbf{Infantil}, precisamos concentrar atenção sobre as condições de aprendizagem dos \textbf{bebês}, \textbf{crianças} bem \textbf{pequenas} e \textbf{crianças} \textbf{pequenas}, como elas se desenvolvem e como o tempo, os espaços e os materiais são organizados para garantir isso.
Esse conjunto de elementos articulados entre si constituem os Direitos de Aprendizagem e Desenvolvimento.
Segundo a BNCC ( 2017, p. 35 ), os direitos “ asseguram, na Educação \textbf{Infantil}, as condições para que as \textbf{crianças} aprendam em situações nas quais possam desempenhar um papel ativo em ambientes que as \textbf{convidem} a vivenciar desafios e a sentirem -se provocadas a resolvê -los, nas quais possam construir significados sobre si, os outros e o mundo social e natural ”.
Os direitos de aprendizagem e desenvolvimento se efetivam na prática \textbf{diária} por meio dos objetivos dos Campos de Experiências que dialogam de forma estreita com as experiências apontadas nas DCNEI a partir dos eixos do currículo Interações e \textbf{Brincadeiras}.
Na Educação \textbf{Infantil} não há uma fragmentação em conteúdos possíveis de divisão em bimestres, trimestres ou semestres.
Os campos de experiências e seus respectivos objetivos de cada grupo \textbf{etário} de \textbf{crianças} devem ser trabalhados ao longo de todo o ano, mantendo viva a preocupação da garantia dos direitos de aprendizagem.
Os \textbf{bebês} e \textbf{crianças} têm direito de Conviver, \textbf{Brincar}, Participar, Explorar, Expressar e Conhecer -se.
Isso nos diz que a organização da proposta de cada \textbf{instituição} não deve privilegiar uma \textbf{criança} deitada ou sentada, fazendo tarefas em folhas de papel e em silêncio.
Se há no espaço institucional e no planejamento de cada professora/professor, no compromisso da garantia dos direitos, teremos os \textbf{bebês} e \textbf{crianças} descritos conforme as DCNEI e BNCC : ativos e capazes de aprender e se desenvolver de forma integral, fazendo uso das variadas linguagens para apreender o mundo e expressar sobre isso.
Em um diálogo afinado com as DCNEI, a BNCC traz o direito de CONVIVER.
É nas interações com os outros que os \textbf{bebês} e as \textbf{crianças} vivenciam a cultura e se relacionam com os seus grupos de convivência, elaborando suas hipóteses sobre o mundo.
Conviver é um direito fundamental para a garantia do desenvolvimento das funções psíquicas superiores, como a linguagem, o pensamento, a memória, a atenção voluntária, o raciocínio, a \textbf{percepção}, o controle da conduta, entre outras funções que caracterizam nossa \textbf{humanização}.
Os \textbf{bebês} e as \textbf{crianças} têm direito de CONVIVER com outras \textbf{crianças} e \textbf{adultos}, em \textbf{pequenos} e grandes grupos, utilizando diferentes linguagens, ampliando o conhecimento de si e do outro, o respeito em relação à cultura e às diferenças entre as pessoas.
E como podemos assegurar esse direito?
Promovendo situações em que o tempo, espaço e materiais sejam organizados para que as relações entre os \textbf{bebês} e as \textbf{crianças}, seus \textbf{pares} e \textbf{adultos} ocorram de forma saudável e com respeito.
As \textbf{crianças} devem ter assegurado o direito de falar e \textbf{ouvir} nas \textbf{rodas} de conversas, durante a contação de história e nas demais propostas da \textbf{rotina}.
Elas têm o que dizer sobre o mundo, mas nem sempre há pessoas \textbf{atentas} e com interesse real sobre essas observações.
É preciso uma \textbf{escuta} sincera e interessada.
Aquilo que a \textbf{criança} fala precisa ser considerado.
Também é necessário promover momentos em que \textbf{crianças} de diferentes idades ampliem contato, \textbf{brinquem} livremente e com jogos de regras, criados por elas juntamente com os \textbf{adultos}.
Diferentes formas de cultura devem ser apresentadas para que as \textbf{crianças} saibam que existem diferentes formas de viver no mundo.
Quando consideramos a \textbf{criança} como uma pessoa que merece respeito e deve ser \textbf{ouvida}, mudamos as relações de convivência que são ampliadoras de um desenvolvimento integral.
O direito de \textbf{BRINCAR} precisa ser assegurado aos \textbf{bebês} e às \textbf{crianças}, dialogando com o eixo do currículo das DCNEI, que nos diz sobre a forma como se aprende e se desenvolve na primeira \textbf{infância}.
Aos \textbf{bebês} deve ser garantidos atenção e tempo dos \textbf{adultos} para que \textbf{conversem}, toquem e apresentem os objetos e seus usos sociais, que logo se tornarão \textbf{interessantes} para eles.
Quando a \textbf{criança} começa a \textbf{imitar} os papéis sociais dos \textbf{adultos} e, posteriormente, caminha para a representação, um importante salto psíquico está ocorrendo em seu desenvolvimento.
A \textbf{brincadeira} de faz de conta ( ou jogo protagonizado, jogo de papéis, \textbf{brincadeira} de papéis ) se torna cada vez mais \textbf{interessante} para a \textbf{criança} porque ela está descobrindo como a sociedade funciona e representa, em suas \textbf{brincadeiras}, essas relações.
O faz de conta é a atividade principal da periodização psíquica da \textbf{criança}, juntamente com as atividades \textbf{plásticas}.
A \textbf{brincadeira} em sua forma mais expandida é essencial para o desenvolvimento da \textbf{criança}.
As \textbf{crianças} têm direito de \textbf{BRINCAR} cotidianamente de diversas formas, em diferentes espaços e tempos, com diferentes \textbf{parceiros} ( \textbf{crianças} e \textbf{adultos} ), ampliando e diversificando seu acesso a produções culturais, seu conhecimento, sua imaginação, sua criatividade, suas experiências emocionais, corporais, \textbf{sensoriais}, expressivas, cognitivas, sociais e \textbf{relacionais}.
E como podemos assegurar esse direito? \textbf{Brincadeiras} livres e dirigidas, em espaços que a promovam, como áreas abertas e os cantinhos com os \textbf{brinquedos} estruturados ( \textbf{brinquedos} industrializados ), construídos ( para \textbf{brincar} de casinha, cozinha, salão de beleza, vendas, profissões etc. ) e não estruturados ( pedaços de madeira, sementes, tampas, \textbf{embalagens}, caixas etc., e objetos do cotidiano, como garrafas, frascos de perfume, \textbf{colher} de pau, cuias, potes, óculos, chapéus, relógios etc., garantido na seleção a segurança da \textbf{criança} ).
A participação do \textbf{adulto} é sempre necessária, seja observando as linguagens e as relações que se desenrolam, seja participando da \textbf{brincadeira}, tornando -a mais complexa e rica, desafiando as hipóteses das \textbf{crianças}.
As \textbf{crianças} não nascem sabendo \textbf{brincar}, elas aprendem a \textbf{brincar}, pois a \textbf{brincadeira} é um ato social e precisa de outros \textbf{parceiros} mais experientes.
PARTICIPAR é um direito assegurado pela BNCC.
Como já dissemos, a \textbf{criança} tem o que falar, porque lida diariamente com o mundo e sobre ele constrói hipóteses.
Ela precisa expressar essas \textbf{percepções} e ter seus horizontes alargados e são os \textbf{adultos} que precisam garantir espaços de interlocuções para que as \textbf{crianças} participem das decisões coletivas tomadas nos espaços institucionais, como as regras de boa convivência.
Elas não devem ser impostas pelos \textbf{adultos}, mas construídas com as \textbf{crianças}, por elas desenhadas e expostas nas \textbf{paredes} de sua sala de referência.
É possível democratizar a participação das \textbf{crianças} em ações cotidianas onde elas sejam \textbf{ouvidas} e discutam formas de resolver situações ou apresentar propostas para determinadas tarefas a serem realizadas.
A atitude dos \textbf{adultos} que conversa com os \textbf{bebês} e explica as ações que serão desenvolvidas \textbf{demonstra} respeito ao direito dos \textbf{bebês} de participarem do seu processo de \textbf{cuidado} e educação.
As \textbf{crianças} têm direito de PARTICIPAR ativamente, com \textbf{adultos} e outras \textbf{crianças}, tanto do planejamento da gestão da escola e das atividades propostas pelo educador quanto da realização das atividades da vida cotidiana, tais como a escolha das \textbf{brincadeiras}, dos materiais e dos ambientes, desenvolvendo diferentes linguagens e elaborando conhecimentos, decidindo e se posicionando.
E como podemos assegurar esse direito?
Envolvendo as \textbf{crianças} nas tomadas de decisões, em que propostas podem ser apresentadas para serem debatidas e escolhidas ou após a discussão do problema ou da necessidade, \textbf{convidar} as \textbf{crianças} a sugerirem propostas.
Um exemplo para colocar esse direito em prática é a ambientação das \textbf{paredes} da sala de referência.
É comum que antes das \textbf{crianças} chegarem, no início do ano, as \textbf{paredes} estejam todas enfeitadas e com cartazes prontos do tempo, da chamada e das regras.
É \textbf{interessante} construir com elas a maior parte desse material.
Isso dará maior sentido para a \textbf{criança}, que se verá como parte importante do processo de ambientação, deixando suas próprias \textbf{marcas}.
Decidir coletivamente as \textbf{brincadeiras}, os \textbf{livros}, às músicas etc., amplia o sentimento de pertencimento e participação da \textbf{criança}.
Os \textbf{bebês} e as \textbf{crianças} percebem o mundo pelos sentidos e eles estão muito sensíveis a esse processo.
Desenvolver os sentidos sociais deve ser nossa preocupação.
Uma formação estética dos sentidos contribui para o processo de \textbf{humanização} de cada \textbf{criança}.
Humanizar -se é aprender a cultura da vida em sociedade e se apropriar daquilo que de mais elevado a humanidade já produziu.
Quando a \textbf{criança} aprende sobre os objetos da cultura, aprende como eles funcionam e para que servem e passa a usá -los com o fim para o qual eles foram criados, está apreendendo séculos de desenvolvimento cultural de sua espécie.
Daí a importância de os espaços institucionais garantirem o direito de EXPLORAR, de sentir o mundo usando todo o seu corpo.
Os \textbf{bebês} e as \textbf{crianças} têm direito de EXPLORAR movimentos, \textbf{gestos}, \textbf{sons}, formas, \textbf{texturas}, cores, palavras, emoções, transformações, relacionamentos, histórias, objetos, elementos da natureza, na escola e fora dela, ampliando seus saberes sobre a cultura, em suas diversas modalidades : as artes, a escrita, a ciência e a tecnologia.
E como podemos assegurar esse direito?
Garantido tempo, espaço e materiais que os \textbf{bebês} e as \textbf{crianças} possam explorar de forma livre.
Eles devem ser disponibilizados de forma que possam usar seu corpo, explorando seus sentidos. \textbf{Bebês} podem ter o cesto de tesouros e tapetes \textbf{sensoriais} construídos nas próprias \textbf{instituições}, sempre observando sua segurança em relação aos objetos.
As \textbf{crianças} bem \textbf{pequenas}, em sua fase de manipulação objetal, precisam tocar, cheirar, sentir com todo o corpo, \textbf{ouvir}, lamber.
Está descobrindo o mundo por meio dos sentidos. \textbf{Crianças} \textbf{pequenas} têm muitas hipóteses sobre as coisas e \textbf{curiosidades} sobre como essas coisas funcionam.
Isso pode se constituir base para planejar momentos de ampliação dessas hipóteses na exploração do seu meio.
Quando o \textbf{bebê} e a \textbf{criança} se apropriam dos objetos da cultura e da cultura imaterial, elas sentem necessidade de se objetivar, de se EXPRESSAR sobre essas \textbf{percepções}.
Vamos percebendo como os direitos dialogam entre si, porque quando garantimos o direito da convivência, da \textbf{brincadeira}, da participação e da exploração, garantimos o direito da expressão.
Essa \textbf{criança} ativa, que tem espaço e voz, tem como se expressar, dizer sobre suas impressões e hipóteses sobre aquilo que ela observa e lhe toca.
Falar e \textbf{ouvir} o outro constitui partes importantes desse processo, mas não são as únicas.
Precisamos garantir outras formas de expressão por meio das múltiplas linguagens, como o desenho, a \textbf{pintura}, a modelagem e a construção.
Os \textbf{bebês} e as \textbf{crianças} têm direito de EXPRESSAR como sujeito dialógico, criativo e sensível, suas necessidades, emoções, sentimentos, dúvidas, hipóteses, descobertas, opiniões, questionamentos, por meio de diferentes linguagens.
E como podemos assegurar esse direito? \textbf{Acolhendo} as necessidades dos \textbf{bebês} e \textbf{crianças}, acalentado seu choro e tentando entender os motivos de suas atitudes.
As \textbf{rodas} de conversa se convertem em momentos cruciais para isso.
As falas das \textbf{crianças} precisam ser incentivadas e \textbf{ouvidas}.
Substituir propostas prontas ( como desenhos para colorir ) por propostas pensadas com as \textbf{crianças} ( como propor diferentes formas de uso do corpo, de suportes e marcadores para desenhar, pintar, modelar e construir – ver discussão sobre múltiplas linguagens ), ampliando a possibilidade de expressão quando usam o desenho como uma linguagem privilegiada para as \textbf{crianças} bem \textbf{pequenas} e \textbf{crianças} \textbf{pequenas}.
Usar a voz para cantar e o corpo para se movimentar também amplia as possibilidades de expressão.
Devemos tomar \textbf{cuidado} com apresentações ensaiadas apenas para \textbf{satisfazer} \textbf{adultos} e que muitas vezes, cansam e geram ansiedade na \textbf{criança}.
Cantar, dançar e \textbf{dramatizar} deve fazer sentido e se constituir uma linguagem expressiva da \textbf{criança}.
A garantia dos cinco direitos já discutidos culmina na garantia do direito dos \textbf{bebês} e das \textbf{crianças} de CONHECER -SE, de conhecer a si e aos outros e se perceber pertencente a um grupo cultural, onde ela também produz cultura.
A atenção dos \textbf{adultos} dirigida aos \textbf{bebês} e às \textbf{crianças} tem um papel especial nessa construção para que eles se percebam um ser único e com características próprias.
Dedicar tempo ao \textbf{cuidado} dispensado aos \textbf{bebês} e às \textbf{crianças} ajuda na compreensão do seu corpo e de suas \textbf{vontades}.
Para a garantia desse direito, precisamos \textbf{acolher} a \textbf{criança} e ajudá -la a entender suas emoções e sentimentos, nomeando -os e explicando suas reações e o que elas representam.
Os \textbf{bebês} e as \textbf{crianças} têm direito de CONHECER -SE e construir sua identidade pessoal, social e cultural, constituindo uma imagem positiva de si e de seus grupos de pertencimento, nas diversas experiências de \textbf{cuidados}, interações, \textbf{brincadeiras} e linguagens vivenciadas na \textbf{instituição} escolar e em seu contexto familiar e comunitário.
E como podemos assegurar esse direito?
Garantindo que os \textbf{bebês} e as \textbf{crianças} se sintam aceitas, \textbf{cuidadas} e protegidas, que sua fala será considerada, assim como seus sentimentos, que não serão ridicularizados ou menosprezados.
É preciso promover experiências de \textbf{cuidado} um do outro.
Quando \textbf{crianças} se desentendem, o caminho que mais promove a reflexão é \textbf{cuidar} daquele que foi machucado para que a outra \textbf{criança} perceba o efeito da sua ação.
Castigar ou deixar a \textbf{criança} isolada para pensar no que fez não contribui para que ela reflita sobre suas ações e pense porque agiu dessa forma.
Ela não \textbf{conseguirá} fazer isso sozinha.
Precisa dos \textbf{adultos} para perceber as consequências de suas ações e as razões pelas quais agiu assim.
Reiterando, precisamos organizar tempo, espaços e materiais que deem conta de promover situações de valorização das atitudes dos \textbf{bebês} e \textbf{crianças} que observam e têm o que falar sobre o mundo, que se relaciona e interage com seus \textbf{pares} e com os \textbf{adultos} ( logo, não apenas com seus professores, mas com todos os funcionários da equipe escolar ), que é \textbf{ouvida} e tem espaço e atenção para suas ideias e sentimentos, que o tempo e as condições da \textbf{brincadeira} sejam um compromisso dos \textbf{adultos} ( e que estes participem da \textbf{brincadeira}, desafiando e enriquecendo as propostas, mas que também deem tempo livre e observe as relações e trocas que as \textbf{crianças} estabelecem entre si ) e que encontra nos \textbf{adultos}, pessoas que lhes auxiliem a compreender suas emoções e vivências.
Estar na Educação \textbf{Infantil} é um grande desafio profissional.
Cada professora/professor deve ser um profundo conhecedor do desenvolvimento dos \textbf{bebês} e \textbf{crianças}.
Isso nos dá condições de planejar fazeres coletivos que incidem sobre o desenvolvimento individual dos \textbf{bebês} e \textbf{crianças}.
Somente compreendendo como elas aprendem e se desenvolvem, poderemos garantir esses direitos, que precisam ser assegurados no \textbf{dia} a \textbf{dia} das \textbf{creches} e \textbf{pré-escolas}.
Os ambientes assim pensados contribuirão para a promoção de desafios e construção de significados diante das experiências vividas pelos \textbf{bebês} e \textbf{crianças}.

% matched lemmas: acolher, adulto, atentar, bebê, brincadeira, brincar, brinquedo, colher, conseguir, conversar, convidar, creche, criança, cuidado, cuidar, curiosidade, demonstrar, dia, diário, dramatizar, embalagem, escuta, etário, gesto, humanização, imitar, infantil, infância, instituição, interessante, livro, marca, ouvir, par, parceiro, parede, pequeno, percepção, pintura, plástico, pré-escola, relacional, roda, rotina, satisfazer, sensorial, som, textura, vontade
\end{textsample}
